% Optional additive fragment. Required: amsmath, amssymb, theorem/proof environments.
% Its labels do not collide with the standalone article's tlbc: prefix.
\subsection{Tilted reciprocal identities and a binary cancellation theorem}
\label{tiltbin:sec:summary}
Write $F_{a,s}(y)=y\Phi(-y,s,a)$ for $y>0$ and $\operatorname{Re}a>0$.
The tilted Lerch continuation \cite{tlbc:report} gives the following ordinary
absolutely convergent integral for every $s,t\in\mathbb C$:
\begin{equation}\label{tiltbin:eq:pair}
 \int_0^\infty y^\delta F_{a+\delta,s}(y)F_{a,t}(1/y)\frac{dy}{y}
 =\frac{\pi}{\sin(\pi\delta)}
       \{\zeta(s+t,a)-\zeta(s+t,a+\delta)\},
\end{equation}
where $|\operatorname{Re}\delta|<1$, $\operatorname{Re}a>0$, and
$\operatorname{Re}(a+\delta)>0$. The difference is formed before the pole
at $s+t=1$ or the tilt collision $\delta=0$ is removed. Its total-order
jets yield generalized Stieltjes differences; its order-zero derivative
yields a Gamma ratio. The full report proves an arbitrary-depth finite
matched-shift master and a normally convergent independent-shift expansion.

For the arithmetic consequence, let $N_1,\ldots,N_r$ be distinct
nonnegative integers and put
\[
 c_j=\prod_{k\ne j}(N_k-N_j)^{-1},\qquad
 S=\sum_jc_j(-1)^{N_j}.
\]
If both parities occur, form
$E(x)=\prod_{N_j\text{ even}}(x-N_j)$ and
$O(x)=\prod_{N_j\text{ odd}}(x-N_j)$, with degrees $e,o$.
Let $U=E^{-1}\bmod O$, normalized by $\deg U<o$.
Then
\begin{equation}\label{tiltbin:eq:binary}
 S=2(-1)^r[x^{o-1}]U,
 \qquad S/2\equiv\binom{r-2}{o-1}\pmod{2\mathbb Z_{(2)}}.
\end{equation}
Here $\mathbb Z_{(2)}$ consists of rational numbers with odd denominator.
Consequently, at $r=2^h+1$ with $h\ge1$, $S=0$ if and only if all shifts
have the same parity. In contrast, for every $k\ge2$ the mixed-parity set
\[
 (0,1,2,2k,2k+1,2k+3,2k+4)
\]
has $S=0$. Its inverse polynomial is
$U(x)=(2x-6k-5)/(3(2k-1)(2k+1)(2k+3))$, so the conclusion follows
from a symbolic degree drop. The associated depth-seven logarithmic
remainder convolutions therefore reduce to rational polynomials in
$\pi^2$.

These identities do not resolve the short Gaussian $S_6$ reduction or
the revised $S_8$ candidate. The mixed-parity family extends the earlier
sufficient equal-parity condition; it does not contradict that condition.
